\documentclass[
  12pt,
  a4paper,
  oneside
]{article}

\usepackage{iftex}
\ifPDFTeX
  \usepackage[utf8]{inputenc}
  \usepackage[T1]{fontenc}
  \usepackage{lmodern}
\else
  \usepackage{fontspec}
  \setmainfont{lmroman12-regular.otf}[
    BoldFont=lmroman12-bold.otf,
    ItalicFont=lmroman12-italic.otf
  ]
  \setmonofont{lmmono12-regular.otf}
\fi
\usepackage[ngerman]{babel}
\usepackage{microtype}

\usepackage[
  left=2.5cm,
  right=2.5cm,
  top=2.5cm,
  bottom=2.5cm
]{geometry}
\usepackage{setspace}
\onehalfspacing

\usepackage{array}
\usepackage{booktabs}
\usepackage{tabularx}
\usepackage{amsmath}
\usepackage{float}

\usepackage{csquotes}
\usepackage[
  backend=biber,
  style=authoryear,
  sorting=nyt,
  giveninits=true,
  maxcitenames=2,
  maxbibnames=99
]{biblatex}
\addbibresource{literatur.bib}

\usepackage{graphicx}
\usepackage[
  colorlinks=true,
  linkcolor=black,
  urlcolor=blue,
  citecolor=black
]{hyperref}

\hypersetup{
  pdftitle={Projektmanagement},
  pdfauthor={Daniel Tempel}
}

\setcounter{secnumdepth}{3}
\setcounter{tocdepth}{3}

% Unsichtbarer Platzhalter, damit der noch leere Gliederungspunkt im PDF erscheint.
\newcommand{\contentplaceholder}{\mbox{}\par}

% Erlaubt bei langen technischen Bezeichnern einen Zeilenumbruch nach dem Inhalt.
\newcommand{\inlinecode}[1]{\texttt{#1}\allowbreak}

\begin{document}

\tableofcontents
\clearpage

% Dieser Auszug ist dem gemeinsamen Kapitel 8 zugeordnet.
\setcounter{section}{8}

\subsection*{Persönliche Reflexion -- Daniel Tempel}

\noindent\textbf{Daniel Tempel}\par
\nopagebreak

Mein Schwerpunkt im Projekt lag auf der Systemarchitektur sowie der
Konzeption und Implementierung des Avisierungsservice. Dazu gehörten
insbesondere die fachliche und technische Modellierung des
Avisierungsprozesses, die Workflow-Steuerung mit Camunda, die
Mandanten- und Sicherheitsmechanismen, das Abfragemodell für das
Avisierungsdashboard sowie die technische Integrationsumgebung und
automatisierte Backend-Tests.

Eine der größten Herausforderungen war die Modellierung des
Avisierungsprozesses als länger laufender Workflow. In einer frühen
Umsetzung wurde ein aus DISPO übermittelter Avisierungsauftrag zunächst
persistiert und die Kundenbenachrichtigung synchron versendet. Camunda
übernahm anschließend im Wesentlichen nur die Überwachung der
Antwortfrist, während die Kundenreaktion außerhalb des Workflows
verarbeitet wurde. Im weiteren Projektverlauf zeigte sich, dass diese
Aufteilung den fachlichen Lebenszyklus nicht ausreichend zusammenhängend
abbildete.

Der Workflow wurde deshalb grundlegend überarbeitet. Der
Nachrichtenversand wurde in die asynchrone Prozessverarbeitung
überführt, während die Workflow Engine die Koordination von Versand,
Kundenreaktion, Fristablauf und Ablösung einer Bestätigungsanfrage
übernahm. Dadurch wurde die Trennung zwischen synchroner
Auftragsübernahme, fachlichem Zustand und technischer Ablaufsteuerung
deutlicher.

Auch die Architektur des Avisierungsservice entwickelte sich während
des Projekts weiter. Nach entsprechendem Feedback wurden die
Verantwortlichkeiten zwischen Domain, Application und Adaptern klarer
getrennt und ausgewählte externe Integrationen über Ports entkoppelt.
Für die Persistenz wurde dagegen bewusst eine pragmatischere Lösung mit
JPA und Spring Data beibehalten. Daraus entstand die Erkenntnis, dass
Architekturprinzipien nicht schematisch angewendet werden sollten,
sondern ihr Nutzen gegen zusätzlichen Struktur- und
Implementierungsaufwand abzuwägen ist.

Die Zusammenarbeit im Team erfolgte überwiegend über GitHub Issues,
getrennte Branches und regelmäßige Abstimmungen. Die Verteilung und
Nachverfolgung der Arbeitspakete wurde von mir in weiten Teilen
koordiniert. Besonders bei Anforderungen, die mehrere
Verantwortungsbereiche betrafen, beispielsweise bei der späteren
Ausgestaltung des Avisierungsdashboards, waren wiederholte fachliche
und technische Abstimmungen erforderlich.

Für mich lagen die wichtigsten Lernergebnisse in der Arbeit mit
länger laufenden und ereignisabhängigen Prozessen, der Abgrenzung von
fachlichen und technischen Zuständen sowie der Behandlung von
Mandantentrennung, Wiederholungen und Nebenläufigkeit. Als zweckmäßig
erwiesen sich insbesondere das Strategy Pattern für die
Kommunikationskanäle, Flyway für die Schemaentwicklung und Docker
Compose für die lokale Integrationsumgebung. Für zukünftige Projekte
würde ich den vollständigen fachlichen Prozess mit seinen Zuständen,
Ereignissen und Verantwortlichkeiten früher modellieren und den
Projektaufwand von Beginn an kontinuierlich arbeitspaketbezogen
erfassen.

\end{document}
